
\chapter{2021年天津卷}

\section{单选题}

\begin{problem}
设集合 \( A=\{-1,0,1\} \)，\( B=\{1,3,5\} \)，\( C=\{0,2,4\} \)，则 \( \left( A\cap B \right)\cup C= \)（\quad）

\choices
    {\(\{0\}\)}
    {\(\{0,1,3,5\}\)}
    {\(\{0,1,2,4\}\)}
    {\(\{0,2,3,4\}\)}

\end{problem}

\begin{answer}
C
\end{answer}

\begin{solution}
因为 \( A\cap B=\{1\} \)，所以 \( \left( A\cap B \right)\cup C=\{0,1,2,4\} \).

故选 C.
\end{solution}

\begin{problem}
已知 \( a\in \R \)，则“\( a>6 \)”是“\( a^2>36 \)”的（\quad）

\choices
    {充分不必要条件}
    {必要不充分条件}
    {充要条件}
    {既不充分也不必要条件}

\end{problem}

\begin{answer}
A
\end{answer}

\begin{solution}
由题意，若 \( a>6 \)，则 \( a^2>36 \)，故充分性成立；

若 \( a^2>36 \)，则 \( a>6 \) 或 \( a<-6 \)，推不出 \( a>6 \)，故必要性不成立；

所以“\( a>6 \)”是“\( a^2>36 \)”的充分不必要条件.

故选 A.
\end{solution}

\begin{problem}
函数 \( y=\dfrac{\ln |x|}{x^2+2} \) 的图像大致为（\quad）

\begin{center}
\bitmapinclude[max width=6.0cm,max totalheight=4.4cm]{img/2021/tianjin/q3_fig1.png}

\bitmapinclude[max width=6.0cm,max totalheight=4.4cm]{img/2021/tianjin/q3_fig2.png}
\end{center}

\end{problem}

\begin{answer}
B
\end{answer}

\begin{solution}
设 \( y=f(x)=\dfrac{\ln |x|}{x^2+2} \)，则函数 \( f(x) \) 的定义域为 \( \{x\mid x\ne 0\} \)，关于原点对称，

又 \( f(-x)=\dfrac{\ln |-x|}{(-x)^2+2}=f(x) \)，所以函数 \( f(x) \) 为偶函数，排除 A、C；

当 \( x\in \left( 0,1 \right) \) 时，\( \ln |x|<0 \)，\( x^2+2>0 \)，所以 \( f(x)<0 \)，排除 D.

故选 B.
\end{solution}

\begin{problem}
从某网络平台推荐的影视作品中抽取 \( 400 \) 部，统计其评分分布数据，将所得 \( 400 \) 个评分数据分为 \( 8 \) 组：
\( [66,70) \)、\( [70,74) \)、\(\cdots\)、\( [94,98] \)，并整理得到如下的频率分布直方图，则评分在区间 \( [82,86) \) 内的影视作品数量是（\quad）

\bitmapfigure[max width=5.5cm,max totalheight=4.7cm]{img/2021/tianjin/q4_fig1.png}

\choices
    {20}
    {40}
    {64}
    {80}

\end{problem}

\begin{answer}
D
\end{answer}

\begin{solution}
由频率分布直方图可知，评分在区间 \( [82,86) \) 内的影视作品数量为
\(400\times 0.05\times 4=80\).

故选 D.
\end{solution}

\begin{problem}
设 \( a=\log_{2}0.3 \)，\( b=\log_{\frac{1}{2}}0.4 \)，\( c=0.4^{0.3} \)，则 \( a,b,c \) 的大小关系为（\quad）

\choices
    {\(a<b<c\)}
    {\(c<a<b\)}
    {\(b<c<a\)}
    {\(a<c<b\)}

\end{problem}

\begin{answer}
D
\end{answer}

\begin{solution}
因为 \( \log_{2}0.3<\log_{2}1=0 \)，所以 \( a<0 \)；

又 \( \log_{\frac{1}{2}}0.4=-\log_{2}0.4=\log_{2}\dfrac{5}{2}>\log_{2}2=1 \)，所以 \( b>1 \)；

因为 \( 0<0.4^{0.3}<0.4^0=1 \)，所以 \( 0<c<1 \).

故 \( a<c<b \).

故选 D.
\end{solution}

\begin{problem}
两个圆锥的底面是一个球的同一截面，顶点均在球面上，若球的体积为 \( \dfrac{32\pi}{3} \)，两个圆锥的高之比为 \( 1:3 \)，则这两个圆锥的体积之和为（\quad）

\choices
    {\(3\pi\)}
    {\(4\pi\)}
    {\(9\pi\)}
    {\(12\pi\)}

\end{problem}

\begin{answer}
B
\end{answer}

\begin{solution}
如下图所示，设两个圆锥的底面圆圆心为点 \( D \)，设圆锥 \( AD \) 和圆锥 \( BD \) 的高之比为 \( 3:1 \)，即 \( AD=3BD \).

\bitmapfigure[max width=3.6cm,max totalheight=4.0cm]{img/2021/tianjin/q6_fig1.png}

设球的半径为 \( R \)，则
\(\frac{4\pi R^3}{3}=\frac{32\pi}{3}\)，
可得 \( R=2 \)，所以 \( AB=AD+BD=4BD=4 \).

所以 \( BD=1 \)，\( AD=3 \).

因为 \( CD\perp AB \)，则 \( \angle CAD+\angle ACD=\angle BCD+\angle ACD=90^\circ \)，所以 \( \angle CAD=\angle BCD \).

又因为 \( \angle ADC=\angle BDC \)，所以 \( \triangle ACD\sim \triangle CBD \)，
\(\frac{AD}{CD}=\frac{CD}{BD}\)，
所以 \( CD=\sqrt{AD\cdot BD}=\sqrt3 \).

因此，这两个圆锥的体积之和为
\(\frac{1}{3}\pi\cdot CD^2\cdot \left( AD+BD \right) =\frac{1}{3}\pi\times 3\times 4 =4\pi\).

故选 B.
\end{solution}

\begin{problem}
若 \( 2^a=5^b=10 \)，则 \( \dfrac{1}{a}+\dfrac{1}{b}= \)（\quad）

\choices
    {\(-1\)}
    {\(\lg 7\)}
    {\(1\)}
    {\(\log_{7}10\)}

\end{problem}

\begin{answer}
C
\end{answer}

\begin{solution}
因为 \( 2^a=5^b=10 \)，所以 \( a=\log_{2}10 \)，\( b=\log_{5}10 \)，所以 \( \frac{1}{a}+\frac{1}{b}=\frac{1}{\log_{2}10}+\frac{1}{\log_{5}10}=\lg 2+\lg 5=\lg 10=1 \).

故选 C.
\end{solution}

\begin{problem}
已知双曲线 \( \dfrac{x^2}{a^2}-\dfrac{y^2}{b^2}=1 \)\( \left( a>0,b>0 \right) \) 的右焦点与抛物线 \( y^2=2px \)\( \left( p>0 \right) \) 的焦点重合，抛物线的准线交双曲线于 \( A \)、\( B \) 两点，交双曲线的渐近线于 \( C \)、\( D \) 两点，若 \( |CD|=\sqrt2\,|AB| \)，则双曲线的离心率为（\quad）

\choices
    {\(\sqrt2\)}
    {\(\sqrt3\)}
    {\(2\)}
    {\(3\)}

\end{problem}

\begin{answer}
A
\end{answer}

\begin{solution}
设双曲线 \( \dfrac{x^2}{a^2}-\dfrac{y^2}{b^2}=1 \)\( \left( a>0,b>0 \right) \) 与抛物线 \( y^2=2px \)\( \left( p>0 \right) \) 的公共焦点为 \( \left( c,0 \right) \)，则抛物线的准线为 \( x=-c \).

令 \( x=-c \)，则
\(\frac{c^2}{a^2}-\frac{y^2}{b^2}=1\)，
解得 \( y=\pm \dfrac{b^2}{a} \)，所以 \( |AB|=\dfrac{2b^2}{a} \).

又因为双曲线的渐近线方程为 \( y=\pm \dfrac{b}{a}x \)，所以 \( |CD|=\dfrac{2bc}{a} \).

因为 \( |CD|=\sqrt2\,|AB| \)，所以
\(\frac{2bc}{a}=\sqrt2\cdot \frac{2b^2}{a}\)，
即 \( c=\sqrt2\,b \)，所以
\(a^2=c^2-b^2=\frac{1}{2}c^2\).

故双曲线的离心率
\(e=\frac{c}{a}=\sqrt2\).

故选 A.
\end{solution}

\begin{problem}
设 \( a\in \R \)，函数 \( f(x)=\cos\left( 2\pi x-2\pi a \right)\)（\(x<a\)），\(f(x)=x^2-2(a+1)x+a^2+5\)（\(x\ge a\)），若 \( f(x) \) 在区间 \( \left( 0,+\infty \right) \) 内恰有 \( 6 \) 个零点，则 \( a \) 的取值范围是（\quad）

\choices
    {\(\left( 2,\frac{9}{4} \right]\cup \left( \frac{5}{2},\frac{11}{4} \right]\)}
    {\(\left( \frac{7}{4},2 \right)\cup \left( \frac{5}{2},\frac{11}{4} \right]\)}
    {\(\left( 2,\frac{9}{4} \right]\cup \left[ \frac{11}{4},3 \right)\)}
    {\(\left( \frac{7}{4},2 \right)\cup \left[ \frac{11}{4},3 \right)\)}

\end{problem}

\begin{answer}
A
\end{answer}

\begin{solution}
因为 \( x^2-2(a+1)x+a^2+5=0 \) 最多有 \( 2 \) 个根，所以 \( \cos\left( 2\pi x-2\pi a \right)=0 \) 至少有 \( 4 \) 个根.

由 \( 2\pi x-2\pi a=\dfrac{\pi}{2}+k\pi \)\( \left( k\in \Z \right) \) 可得
\(x=\frac{k}{2}+\frac{1}{4}+a\)，\(k\in \Z\).
由 \( 0<\dfrac{k}{2}+\dfrac{1}{4}+a<a \) 可得
\(-2a-\frac{1}{2}<k<-\frac{1}{2}\).

（1）当 \( x<a \) 时，

当 \( -5\le -2a-\dfrac{1}{2}<-4 \) 时，\( f(x) \) 有 \( 4 \) 个零点，即 \( \dfrac{7}{4}<a\le \dfrac{9}{4} \)；

当 \( -6\le -2a-\dfrac{1}{2}<-5 \) 时，\( f(x) \) 有 \( 5 \) 个零点，即 \( \dfrac{9}{4}<a\le \dfrac{11}{4} \)；

当 \( -7\le -2a-\dfrac{1}{2}<-6 \) 时，\( f(x) \) 有 \( 6 \) 个零点，即 \( \dfrac{11}{4}<a\le \dfrac{13}{4} \).

（2）当 \( x\ge a \) 时，\( f(x)=x^2-2(a+1)x+a^2+5 \)，
\(\Delta=4(a+1)^2-4\left( a^2+5 \right)=8(a-2)\).

当 \( a<2 \) 时，\( \Delta<0 \)，\( f(x) \) 无零点；

当 \( a=2 \) 时，\( \Delta=0 \)，\( f(x) \) 有 \( 1 \) 个零点；

当 \( a>2 \) 时，
\(f(a)=a^2-2a(a+1)+a^2+5=-2a+5\).
若 \( 2<a\le \dfrac{5}{2} \)，则 \( f(x) \) 有 \( 2 \) 个零点；

若 \( a>\dfrac{5}{2} \)，则 \( f(x) \) 有 \( 1 \) 个零点.

综上，要使 \( f(x) \) 在区间 \( \left( 0,+\infty \right) \) 内恰有 \( 6 \) 个零点，则
\(a\in \left( 2,\frac{9}{4} \right]\cup \left( \frac{5}{2},\frac{11}{4} \right]\).

故选 A.
\end{solution}
\section{填空题}

\begin{problem}
\( \i \) 是虚数单位，复数 \( \dfrac{9+2\i}{2+\i}=\fillinblank{} \).

\end{problem}

\begin{answer}
\(4-\i\)
\end{answer}

\begin{solution}
\(\frac{9+2\i}{2+\i} =\frac{(9+2\i)(2-\i)}{(2+\i)(2-\i)} =\frac{20-5\i}{5} =4-\i\).

故答案为：\( 4-\i \).
\end{solution}

\begin{problem}
在 \( \left( 2x^3+\dfrac{1}{x} \right)^6 \) 的展开式中，\( x^6 \) 的系数是 \(\fillinblank{}\).

\end{problem}

\begin{answer}
160
\end{answer}

\begin{solution}
 \( \left( 2x^3+\dfrac{1}{x} \right)^6 \) 的展开式的通项为
\(T_{r+1}=\mathrm{C}_6^r\left( 2x^3 \right)^{6-r}\left( \frac{1}{x} \right)^r =2^{6-r}\mathrm{C}_6^r x^{18-4r}\).

令 \( 18-4r=6 \)，解得 \( r=3 \).

所以 \( x^6 \) 的系数是
\(2^3\mathrm{C}_6^3=160\).

故答案为：160.
\end{solution}

\begin{problem}
若斜率为 \( \sqrt3 \) 的直线与 \( y \) 轴交于点 \( A \)，与圆 \( x^2+\left( y-1 \right)^2=1 \) 相切于点 \( B \)，则 \( |AB|=\fillinblank{} \).

\end{problem}

\begin{answer}
\(\sqrt3\)
\end{answer}

\begin{solution}
设直线 \( AB \) 的方程为 \( y=\sqrt3\,x+b \)，则点 \( A\left( 0,b \right) \).

由于直线 \( AB \) 与圆 \( x^2+\left( y-1 \right)^2=1 \) 相切，且圆心为 \( C\left( 0,1 \right) \)，半径为 \( 1 \)，所以
\(\frac{|b-1|}{2}=1\)，
解得 \( b=-1 \) 或 \( b=3 \)，所以 \( |AC|=2 \).

因为 \( |BC|=1 \)，故
\(|AB|=\sqrt{|AC|^2-|BC|^2}=\sqrt3\).

故答案为：\( \sqrt3 \).
\end{solution}

\begin{problem}
若 \( a>0 \)，\( b>0 \)，则 \( \dfrac{1}{a}+\dfrac{a}{b^2}+b \) 的最小值为 \(\fillinblank{}\).

\end{problem}

\begin{answer}
\(2\sqrt2\)
\end{answer}

\begin{solution}
因为 \( a>0 \)，\( b>0 \)，所以
\(\frac{1}{a}+\frac{a}{b^2}+b \ge 2\sqrt{\frac{1}{a}\cdot \frac{a}{b^2}}+b =\frac{2}{b}+b \ge 2\sqrt2\).

当且仅当
\(\frac{1}{a}=\frac{a}{b^2}\)且\(\frac{2}{b}=b\)，
即 \( a=b=\sqrt2 \) 时等号成立.

故答案为：\( 2\sqrt2 \).
\end{solution}

\begin{problem}
甲、乙两人在每次猜谜活动中各猜一个谜语，若一方猜对且另一方猜错，则猜对的一方获胜，否则本次平局. 已知每次活动中，甲、乙猜对的概率分别为 \( \dfrac{5}{6} \)、\( \dfrac{1}{5} \)，且每次活动中甲、乙猜对与否互不影响，各次活动也互不影响，则一次活动中，甲获胜的概率为 \(\fillinblank{}\)，\( 3 \) 次活动中，甲至少获胜 \( 2 \) 次的概率为 \(\fillinblank{}\).

\end{problem}

\begin{answer}
\(\dfrac{2}{3},\ \dfrac{20}{27}\)
\end{answer}

\begin{solution}
由题可得一次活动中，甲获胜的概率为
\(\frac{5}{6}\times \frac{4}{5}=\frac{2}{3}\).

则在 \( 3 \) 次活动中，甲至少获胜 \( 2 \) 次的概率为
\(\mathrm{C}_3^2\left( \frac{2}{3} \right)^2\times \frac{1}{3}+\left( \frac{2}{3} \right)^3=\frac{20}{27}\).

故答案为：\( \dfrac{2}{3} \)；\( \dfrac{20}{27} \).
\end{solution}

\begin{problem}
在边长为 \( 1 \) 的等边三角形 \( ABC \) 中，\( D \) 为线段 \( BC \) 上的动点，\( DE\perp AB \) 且交 \( AB \) 于点 \( E \)，\( DF//AB \) 且交 \( AC \) 于点 \( F \)，则 \( \left| 2\overrightarrow{BE}+\overrightarrow{DF} \right| \) 的值为 \(\fillinblank{}\)；\( \left( \overrightarrow{DE}+\overrightarrow{DF} \right)\cdot \overrightarrow{DA} \) 的最小值为 \(\fillinblank{}\).

\bitmapfigure[max width=6.0cm,max totalheight=3.6cm]{img/2021/tianjin/q15_fig1.png}

\end{problem}

\begin{answer}
\(1,\ \dfrac{11}{20}\)
\end{answer}

\begin{solution}
设 \( BE=x \)，\( x\in \left( 0,\dfrac{1}{2} \right) \). 因为 \( ABC \) 为边长为 \( 1 \) 的等边三角形，\( DE\perp AB \)，所以
\(\angle BDE=30^\circ\)，\(BD=2x\)，\(DE=\sqrt3\,x\)，\(DC=1-2x\).

因为 \( DF//AB \)，所以 \( DFC \) 为边长为 \( 1-2x \) 的等边三角形，且 \( DE\perp DF \).

故
\(\left( 2\overrightarrow{BE}+\overrightarrow{DF} \right)^2 =4BE^2+4BE\cdot DF+DF^2 =4x^2+4x(1-2x)\cos 0^\circ+\left( 1-2x \right)^2 =1\)，
所以
\(\left| 2\overrightarrow{BE}+\overrightarrow{DF} \right|=1\).

又
又 \( \left( \overrightarrow{DE}+\overrightarrow{DF} \right)\cdot \overrightarrow{DA}=\left( \overrightarrow{DE}+\overrightarrow{DF} \right)\cdot \left( \overrightarrow{DE}+\overrightarrow{EA} \right)=DE^2+DF\cdot EA=\left( \sqrt3\,x \right)^2+\left( 1-2x \right)\left( 1-x \right) =5x^2-3x+1 =5\left( x-\frac{3}{10} \right)^2+\frac{11}{20} \).

所以当 \( x=\dfrac{3}{10} \) 时，\( \left( \overrightarrow{DE}+\overrightarrow{DF} \right)\cdot \overrightarrow{DA} \) 的最小值为 \( \dfrac{11}{20} \).

故答案为：\( 1 \)；\( \dfrac{11}{20} \).
\end{solution}
\section{解答题}

\begin{problem}
在 \( \triangle ABC \) 中，角 \( A \)、\( B \)、\( C \) 所对的边分别为 \( a \)、\( b \)、\( c \)，已知 \( \sin A:\sin B:\sin C=2:1:\sqrt2 \)，\( b=\sqrt2 \).

\begin{enumerate}
\item 求 \( a \) 的值；
\item 求 \( \cos C \) 的值；
\item 求 \( \sin\left( 2C-\dfrac{\pi}{6} \right) \) 的值.
\end{enumerate}

\end{problem}

\begin{solution}
（1）因为 \( \sin A:\sin B:\sin C=2:1:\sqrt2 \)，由正弦定理可得
\(a:b:c=2:1:\sqrt2\).

又 \( b=\sqrt2 \)，所以
\(a=2\sqrt2\)，\(c=2\).

（2）由余弦定理可得
\(\cos C=\frac{a^2+b^2-c^2}{2ab} =\frac{8+2-4}{2\times 2\sqrt2\times \sqrt2} =\frac{3}{4}\).

（3）因为 \( \cos C=\dfrac{3}{4} \)，所以
\(\sin C=\sqrt{1-\cos^2 C}=\frac{\sqrt7}{4}\).

于是
\(\sin 2C=2\sin C\cos C =2\times \frac{\sqrt7}{4}\times \frac{3}{4} =\frac{3\sqrt7}{8}\)，
\(\cos 2C=2\cos^2 C-1 =2\times \frac{9}{16}-1 =\frac{1}{8}\).

故
故 \( \sin\left( 2C-\frac{\pi}{6} \right)=\sin 2C\cos \frac{\pi}{6}-\cos 2C\sin \frac{\pi}{6}=\frac{3\sqrt7}{8}\times \frac{\sqrt3}{2}-\frac{1}{8}\times \frac{1}{2}=\frac{3\sqrt{21}-1}{16} \).
\end{solution}

\begin{problem}
如图，在棱长为 \( 2 \) 的正方体 \( ABCD-A_1B_1C_1D_1 \) 中，\( E \) 为棱 \( BC \) 的中点，\( F \) 为棱 \( CD \) 的中点.

\begin{enumerate}
\item 求证：\( D_1F // \) 平面 \( A_1EC_1 \)；
\item 求直线 \( AC_1 \) 与平面 \( A_1EC_1 \) 所成角的正弦值；
\item 求二面角 \( A-A_1C_1-E \) 的正弦值.
\end{enumerate}

\bitmapfigure[max width=5cm,max totalheight=4.7cm]{img/2021/tianjin/q17_fig1.png}

\end{problem}

\begin{solution}
（1）以 \( A \) 为原点，\( AB \)、\( AD \)、\( AA_1 \) 分别为 \( x \)、\( y \)、\( z \) 轴，建立如图空间直角坐标系.

\bitmapfigure[max width=5cm,max totalheight=4.7cm]{img/2021/tianjin/q17_solution_fig1.png}

则
\(A(0,0,0)\)，\(A_1(0,0,2)\)，\(B(2,0,0)\)，\(C(2,2,0)\)，\(D(0,2,0)\)，
\(C_1(2,2,2)\)，\(D_1(0,2,2)\).

因为 \( E \) 为棱 \( BC \) 的中点，\( F \) 为棱 \( CD \) 的中点，所以
\(E(2,1,0)\)，\(F(1,2,0)\).

于是
\(\overrightarrow{D_1F}=(1,0,-2)\)，\(\overrightarrow{A_1C_1}=(2,2,0)\)，\(\overrightarrow{A_1E}=(2,1,-2)\).

设平面 \( A_1EC_1 \) 的一个法向量为 \( \bs{m}=(x_1,y_1,z_1) \)，则
\examdisplaycases{}{
\bs{m}\cdot \overrightarrow{A_1C_1}=2x_1+2y_1=0,\\
\bs{m}\cdot \overrightarrow{A_1E}=2x_1+y_1-2z_1=0.
}
令 \( x_1=2 \)，可得 \( \bs{m}=(2,-2,1) \).

因为
\(\overrightarrow{D_1F}\cdot \bs{m}=2-2=0\)，
所以 \( D_1F\perp \bs{m} \).

又因为 \( D_1F \not\subset \) 平面 \( A_1EC_1 \)，所以 \( D_1F // \) 平面 \( A_1EC_1 \).

（2）由（1）得
\(\overrightarrow{AC_1}=(2,2,2)\).

设直线 \( AC_1 \) 与平面 \( A_1EC_1 \) 所成角为 \( \theta \)，则 \(\sin \theta=\left| \cos \left\langle \bs{m},\overrightarrow{AC_1} \right\rangle \right|=\frac{|\bs{m}\cdot \overrightarrow{AC_1}|}{|\bs{m}|\,|\overrightarrow{AC_1}|}=\frac{2}{3\times 2\sqrt3}=\frac{\sqrt3}{9}\).

（3）由正方体的特征可得，平面 \( AA_1C_1 \) 的一个法向量为
\(\overrightarrow{DB}=(2,-2,0)\).

于是 \(\left| \cos \left\langle \overrightarrow{DB},\bs{m} \right\rangle \right|=\frac{|\overrightarrow{DB}\cdot \bs{m}|}{|\overrightarrow{DB}|\,|\bs{m}|}=\frac{8}{3\times 2\sqrt2}=\frac{2\sqrt2}{3}\).

所以二面角 \( A-A_1C_1-E \) 的正弦值为
\(\sqrt{1-\cos^2 \left\langle \overrightarrow{DB},\bs{m} \right\rangle }=\frac{1}{3}\).
\end{solution}

\begin{problem}
已知椭圆 \( \dfrac{x^2}{a^2}+\dfrac{y^2}{b^2}=1 \)\( \left( a>b>0 \right) \) 的右焦点为 \( F \)，上顶点为 \( B \)，离心率为 \( \dfrac{2\sqrt5}{5} \)，且 \( |BF|=\sqrt5 \).

\begin{enumerate}
\item 求椭圆的方程；
\item 直线 \( l \) 与椭圆有唯一的公共点 \( M \)，与 \( y \) 轴的正半轴交于点 \( N \)，过 \( N \) 与 \( BF \) 垂直的直线交 \( x \) 轴于点 \( P \). 若 \( MP//BF \)，求直线 \( l \) 的方程.
\end{enumerate}

\end{problem}

\begin{solution}
（1）由椭圆 \( \dfrac{x^2}{a^2}+\dfrac{y^2}{b^2}=1 \)\( \left( a>b>0 \right) \) 可知，其右焦点在 \( x \) 轴上，坐标为 \( F(c,0) \)，上顶点在 \( y \) 轴上，坐标为 \( B(0,b) \)，故
\(|BF|=\sqrt{c^2+b^2}=a=\sqrt5\).

因为椭圆的离心率为
\(e=\frac{c}{a}=\frac{2\sqrt5}{5}\)，
故 \( c=2 \)，\( b=\sqrt{a^2-c^2}=1 \).

因此，椭圆的方程为
\(\frac{x^2}{5}+y^2=1\).

（2）设点 \( M(x_0,y_0) \) 为椭圆 \( \dfrac{x^2}{5}+y^2=1 \) 上一点.

先证明直线 \( MN \) 的方程为 \(\frac{x_0x}{5}+y_0y=1\).

联立
\examdisplaycases{}{
\dfrac{x_0x}{5}+y_0y=1,\\[4pt]
\dfrac{x^2}{5}+y^2=1,
}
消去 \( y \) 并整理得
\(x^2-2x_0x+x_0^2=0,\)
其判别式
\(\Delta=4x_0^2-4x_0^2=0.\)

因此，椭圆 \( \dfrac{x^2}{5}+y^2=1 \) 在点 \( M(x_0,y_0) \) 处的切线方程为 \(\frac{x_0x}{5}+y_0y=1.\)

\bitmapfigure[max width=4.2cm,max totalheight=3.8cm]{img/2021/tianjin/q18_solution_fig1.png}

在直线 \( MN \) 的方程中，令 \( x=0 \)，可得 \( y=\dfrac{1}{y_0} \). 由题意知 \( y_0>0 \)，即点
\(N\left( 0,\frac{1}{y_0} \right)\).

直线 \( BF \) 的斜率为 \(k_{BF}=-\frac{b}{c}=-\frac{1}{2},\) 所以，直线 \( PN \) 的方程为 \(y=2x+\frac{1}{y_0}.\)

在直线 \( PN \) 的方程中，令 \( y=0 \)，可得 \( x=-\dfrac{1}{2y_0} \)，即点
\(P\left( -\frac{1}{2y_0},0 \right)\).

因为 \( MP//BF \)，则 \( k_{MP}=k_{BF} \)，即 \(\frac{y_0}{x_0+\frac{1}{2y_0}}=-\frac{1}{2},\) 整理可得 \(\left( x_0+5y_0 \right)^2=0.\)

所以 \( x_0=-5y_0 \). 又因为 \(\frac{x_0^2}{5}+y_0^2=1,\) 所以 \(6y_0^2=1.\)

由 \( y_0>0 \) 可得
\(y_0=\frac{\sqrt6}{6},\quad x_0=-\frac{5\sqrt6}{6}.\)

故直线 \( l \) 的方程为 \(-\frac{\sqrt6}{6}x+\frac{\sqrt6}{6}y=1,\) 即 \(x-y+\sqrt6=0.\)
\end{solution}

\begin{problem}
已知 \( \{a_n\} \) 是公差为 \( 2 \) 的等差数列，其前 \( 8 \) 项和为 \( 64 \). \( \{b_n\} \) 是公比大于 \( 0 \) 的等比数列，\( b_1=4 \)，\( b_3-b_2=48 \).

\begin{enumerate}
\item 求 \( \{a_n\} \) 和 \( \{b_n\} \) 的通项公式；
\item 记 \( c_n=b_{2n}+\dfrac{1}{b_n} \)，\( n\in \N^* \)，
\begin{enumerate}
\item 证明 \( \{c_n^2-c_{2n}\} \) 是等比数列；
\item 证明 \( \sum_{k=1}^n \sqrt{\dfrac{a_ka_{k+1}}{c_k^2-c_{2k}}}<2\sqrt2 \)\( \left( n\in \N^* \right) \).
\end{enumerate}
\end{enumerate}

\end{problem}

\begin{solution}
（1）因为 \( \{a_n\} \) 是公差为 \( 2 \) 的等差数列，其前 \( 8 \) 项和为 \( 64 \)，所以
\(a_1+a_2+\cdots +a_8=8a_1+\frac{8\times 7}{2}\times 2=64,\)
所以 \( a_1=1 \).

故
\(a_n=a_1+2(n-1)=2n-1,\quad n\in \N^*.\)

设等比数列 \( \{b_n\} \) 的公比为 \( q \)\( \left( q>0 \right) \)，则
\(b_3-b_2=b_1q^2-b_1q=4(q^2-q)=48,\)
解得 \( q=4 \)（负值舍去）.

故
\(b_n=b_1q^{n-1}=4^n,\quad n\in \N^*.\)

（2）（i）由题意，
\(c_n=b_{2n}+\frac{1}{b_n}=4^{2n}+\frac{1}{4^n}.\)

所以
\(c_n^2-c_{2n} =\left( 4^{2n}+\frac{1}{4^n} \right)^2-\left( 4^{4n}+\frac{1}{4^{2n}} \right) =2\cdot 4^n.\)

因此 \( c_n^2-c_{2n}\ne 0 \)，且
\(\frac{c_{n+1}^2-c_{2n+2}}{c_n^2-c_{2n}} =\frac{2\cdot 4^{n+1}}{2\cdot 4^n} =4.\)

所以数列 \( \{c_n^2-c_{2n}\} \) 是等比数列.

（2）由题意知，
由题意知，\( \frac{a_na_{n+1}}{c_n^2-c_{2n}}=\frac{(2n-1)(2n+1)}{2\cdot 4^n}=\frac{4n^2-1}{2\cdot 4^n}<\frac{4n^2}{2\cdot 2^{2n}}=\frac{n^2}{2^{2n-1}} \).

所以
\(\sqrt{\frac{a_na_{n+1}}{c_n^2-c_{2n}}} <\sqrt{\frac{n^2}{2^{2n-1}}} =\frac{n}{\sqrt2\cdot 2^{n-1}}.\)

故
\(\sum_{k=1}^n \sqrt{\frac{a_ka_{k+1}}{c_k^2-c_{2k}}} <\frac{1}{\sqrt2}\sum_{k=1}^n \frac{k}{2^{k-1}}.\)

设
\(T_n=\sum_{k=1}^n \frac{k}{2^{k-1}}=\frac{1}{2^0}+\frac{2}{2^1}+\frac{3}{2^2}+\cdots +\frac{n}{2^{n-1}}.\)

则
\(\frac{1}{2}T_n=\frac{1}{2^1}+\frac{2}{2^2}+\frac{3}{2^3}+\cdots +\frac{n}{2^n}.\)

两式相减得
两式相减得 \( \frac{1}{2}T_n=1+\frac{1}{2^1}+\frac{1}{2^2}+\cdots +\frac{1}{2^{n-1}}-\frac{n}{2^n}=\frac{1\cdot \left( 1-\frac{1}{2^n} \right)}{1-\frac{1}{2}}-\frac{n}{2^n}=\frac{2^n-n-2}{2^n} \).

所以
\(T_n=4-\frac{n+2}{2^{n-1}}.\)

故
\(\sum_{k=1}^n \sqrt{\frac{a_ka_{k+1}}{c_k^2-c_{2k}}} <\frac{1}{\sqrt2}\left( 4-\frac{n+2}{2^{n-1}} \right) <2\sqrt2.\)
\end{solution}

\begin{problem}
已知 \( a>0 \)，函数 \( f(x)=ax-x\e^x \).

\begin{enumerate}
\item 求曲线 \( y=f(x) \) 在点 \( \left( 0,f(0) \right) \) 处的切线方程；
\item 证明 \( f(x) \) 存在唯一的极值点；
\item 若存在 \( a \)，使得 \( f(x)\le a+b \) 对任意 \( x\in \R \) 成立，求实数 \( b \) 的取值范围.
\end{enumerate}

\end{problem}

\begin{solution}
（1）\(f'(x)=a-(x+1)\e^x,\)
则 \( f'(0)=a-1 \)，又 \( f(0)=0 \)，故切线方程为
\(y=(a-1)x,\quad a>0.\)

（2）令
\(f'(x)=a-(x+1)\e^x=0,\)
则
\(a=(x+1)\e^x.\)

令 \( g(x)=(x+1)\e^x \)，则
\(g'(x)=(x+2)\e^x.\)

当 \( x\in \left( -\infty,-2 \right) \) 时，\( g'(x)<0 \)，\( g(x) \) 单调递减；当 \( x\in \left( -2,+\infty \right) \) 时，\( g'(x)>0 \)，\( g(x) \) 单调递增.

当 \( x\to -\infty \) 时，\( g(x)<0 \)，且 \( g(-1)=0 \)；当 \( x\to +\infty \) 时，\( g(x)>0 \). 画出 \( g(x) \) 大致图像如下：

\bitmapfigure[max width=4.6cm,max totalheight=4.0cm]{img/2021/tianjin/q20_solution_fig1.png}

所以当 \( a>0 \) 时，\( y=a \) 与 \( y=g(x) \) 仅有一个交点. 设 \( g(m)=a \)，则 \( m>-1 \)，且
\(f'(m)=a-g(m)=0.\)

当 \( x\in \left( -\infty,m \right) \) 时，\( a>g(x) \)，则 \( f'(x)>0 \)，\( f(x) \) 单调递增；

当 \( x\in \left( m,+\infty \right) \) 时，\( a<g(x) \)，则 \( f'(x)<0 \)，\( f(x) \) 单调递减.

故 \( x=m \) 为 \( f(x) \) 的极大值点，所以 \( f(x) \) 存在唯一的极值点.

（3）由（2）知 \( f(x)_{\max}=f(m) \)，此时 \( a=(1+m)\e^m \)，\( m>-1 \).

所以
\(\{f(x)-a\}_{\max}=f(m)-a=(m^2-m-1)\e^m,\quad m>-1.\)

令
\(h(x)=(x^2-x-1)\e^x,\quad x>-1.\)

若存在 \( a \)，使得 \( f(x)\le a+b \) 对任意 \( x\in \R \) 成立，等价于存在 \( x\in \left( -1,+\infty \right) \)，使得 \( h(x)\le b \)，即 \( b\ge h(x)_{\min} \).

又
\(h'(x)=(x^2+x-2)\e^x=(x-1)(x+2)\e^x,\quad x>-1.\)

当 \( x\in \left( -1,1 \right) \) 时，\( h'(x)<0 \)，\( h(x) \) 单调递减；

当 \( x\in \left( 1,+\infty \right) \) 时，\( h'(x)>0 \)，\( h(x) \) 单调递增.

所以
\(h(x)_{\min}=h(1)=-\e,\)
故 \( b\ge -\e \).

因此实数 \( b \) 的取值范围为
\(\left[ -\e,+\infty \right).\)
\end{solution}
